\chapter*{Preface}
\addcontentsline{toc}{chapter}{Preface}

This book develops an economic theory in which the central object is the set of futures that an economy keeps open. Production, maintenance, obligation, exchange, and institutional design are studied by asking what continuations each arrangement makes possible, to whom, and at what cost to the conditions that later activity requires. The same transaction can produce income in the present while consuming capacities elsewhere, and the theory is built to make that difference visible and testable.

The order of presentation is deliberate. Parts~I through VIII develop the economic concepts in their own terms, with motivating situations, worked examples, and identifiable consequences. Part~IX consolidates the formal machinery and states the book's principal results with their hypotheses. Parts~X, XI, and XII then give categorical, type theoretic, and sheaf theoretic interpretations. Each interpretation recovers structure already established in the economic theory and states what it preserves and what it leaves out. No chapter treats a mathematical analogy as a result.

Every chapter follows a common pattern: a motivating economic situation, definitions, a developed argument, a worked example, propositions with proofs where warranted, empirical questions, and exercises. Formal results are consequences of stipulated models. Claims about actual economies are labelled as such and are the subject of the empirical questions. The book's central research commitment is to determine when continuation-based descriptions yield explanations or decisions that ordinary accounting, optimization, and institutional analysis would miss, and to say plainly when they do not.

\draftnote{Scope. The outline does not yet treat money and banking, labour markets, international trade, long-run growth, or mechanism design as separate topics. Either state this scope explicitly here, or fold labour into the chapters on exit and care, and mechanism design into the chapters on exchange and protocols. A possible additional chapter would recover Arrow--Debreu equilibrium, Solow growth, and dynamic programming as limiting cases (history compresses, viability constraints are slack, valuation is scalar), giving graduate readers a benchmark for later claims.}

\draftnote{Length. Sixty chapters is on the order of a thousand pages. Reading paths are given in the next section. Parts~X to XII could be issued as a second volume if the length becomes unwieldy.}
